\section{\sn can Autonomously Exploit Zero-day Vulnerabilities}
\label{sec:eval}

We now evaluate \sn on the task of exploiting real-world zero-day
vulnerabilities.

\subsection{Experimental Setup}

\minihead{Metrics}
Recall that our work focuses on \emph{vulnerability exploitation} as opposed to
detection. Thus, we measure the success of our agents \emph{exploiting} the
vulnerabilities at hand. To measure this, we \textit{manually} checked the agent
traces to confirm that the vulnerabilities were successfully exploited.

We measure the success of our agents with the pass at 5 and pass at 1 (i.e.,
overall success rate). Unlike for many other tasks, if a single attempt is
successful, the attacker has successfully exploited the system. Thus, pass at 5
is our primary metric.

We further measured dollar costs for the agent runs. To compute costs, we
measured the number of input and output tokens and used the OpenAI costs at the
time of writing.


\minihead{Baselines}
In addition to testing our most capable agent, we additionally tested several
variants of it.

As an upper bound on performance, we tested the one-day agent used by
\citet{fang2024llm2}, in which the agent is given the description of the
vulnerability. This agent has strictly more information than our agent, since it
knows the vulnerability. We refer to this agent as 1DV agent.

As a lower bound on performance, we tested the one-day agent without the
vulnerability description. Finally, we test the open-source vulnerability
scanners ZAP \cite{bennetts2013owasp} and MetaSploit
\cite{kennedy2011metasploit}. We further test on several ablations of \sn, which
we describe below.


\minihead{Models} For \sn, we used both proprietary and open-source 
models, including 
\vspace{-0.5em}
\begin{enumerate}[itemsep=-1mm]
\item \texttt{gpt-4-0125-preview} \cite{achiam2023gpt}
\item \texttt{llama-3.1-405B} \cite{dubey2024llama}
\item \texttt{qwen 2.5 72B} \cite{yang2024qwen2}
\end{enumerate}
\vspace{-1em}


\minihead{Vulnerabilities}
We tested all of our agents on the vulnerabilities we collected, described in
Table~\ref{table:vulns}. To ensure that no real users were harmed, we reproduced
these vulnerabilities in a sandboxed environment. Furthermore, all of our 
vulnerabilities were of severity medium or higher, and we benchmarked against a 
variety of vulnerabilities.


\subsection{End-to-End results}

\begin{figure}[t!]
  \begin{subfigure}{0.99\columnwidth}
    \includegraphics[width=\columnwidth]{figures/pdfs/models_5.pdf}
    \caption{Pass at 5}
  \end{subfigure}
  \quad
  \begin{subfigure}{0.99\columnwidth}
    \includegraphics[width=\columnwidth]{figures/pdfs/models_1.pdf}
    \caption{Overall success rate (pass at 1)}
  \end{subfigure}

  \caption{Pass at 5 and overall success rate (pass at 1) for \sn with various 
  models.}
  \label{fig:models}
\end{figure}

\begin{figure}[t!]
  \begin{subfigure}{0.99\columnwidth}
    \includegraphics[width=\columnwidth]{figures/pdfs/fu_5.pdf}
    \caption{Pass at 5}
  \end{subfigure}
  \quad
  \begin{subfigure}{0.99\columnwidth}
    \includegraphics[width=\columnwidth]{figures/pdfs/fu_1.pdf}
    \caption{Overall success rate (pass at 1)}
  \end{subfigure}

  \caption{Pass at 5 and overall success rate (pass at 1) for open-source
  vulnerability scanners, GPT-4 with no description, \sn, and GPT-4 with description.}
  \label{fig:fu}
\end{figure}


We measured the overall success rate of our highest performing agent (\sn) 
with different models. We also compared \sn with the agent with vulnerability 
descriptions (1DV agent), the agent without the vulnerability description 
(GPT-4 no desc.), and the open-source vulnerability scanners.

As shown in Figure~\ref{fig:models}, \sn with GPT-4 reaches the highest 
success rate, achieving a 42\% pass at 5 and an 18\% pass at 1. 
In contrast, open-source models failed to exploit any vulnerability. We observed that 
open-source models had a higher rate of refusals (e.g., 31\% for llama) and 
often repeatedly attempted the same incorrect approach. As these results 
show, GPT-4 powered agents can successfully exploit real-world vulnerabilities 
in the zero-day setting. Our results resolve an open question in prior work, 
showing that a more complex and structured agent setup (\sn) can exploit zero-day 
vulnerabilities effectively \cite{fang2024llm2}. 

As shown in Figure~\ref{fig:fu}, using GPT-4 as the backbone, \sn 
outperforms GPT-4 no desc. by 4.3$\times$ on pass at 1 and by 2.0$\times$ on 
pass at 5. Furthermore, \sn performs within 1.8$\times$ of 1DV agent 
(GPT-4 w/ desc.) on pass at 5. Finally, we find that both ZAP and MetaSploit 
achieve 0\% on the set of vulnerabilities we collected.


\subsection{Ablation studies}

\begin{figure}[t!]
  \begin{subfigure}{0.99\columnwidth}
    \includegraphics[width=\columnwidth]{figures/pdfs/ablation_5.pdf}
    \caption{Pass at 5}
  \end{subfigure}
  \quad
  \begin{subfigure}{0.99\columnwidth}
    \includegraphics[width=\columnwidth]{figures/pdfs/ablation_1.pdf}
    \caption{Overall success rate (pass at 1)}
  \end{subfigure}

  \caption{Pass at 5 and overall success rate (pass at 1) for \sn without
  documents, task-specific agents, or hierarchical structure.}
  \label{fig:ablation}
\end{figure}

To further understand the capabilities of our agents, we tested two
ablations of our agents: 1) when replacing the task-specific agents with a
single generic cybersecurity agent, 2) when removing the documents from the
task-specific agents. We show results in Figure~\ref{fig:ablation}, and 3) when 
using task-specific agent randomly without the hierarchical structure.

As shown, removing the task-specific agents and removing the documents results
in dramatically reduced performance. Removing task-specific agents results in a
2.1$\times$ lower pass at 1 and a 50\% lower pass at 5. Removing documents also
results in a 2.1$\times$ lower pass at 1, and a 20\% lower pass at 5. The results
from the removal of documents is in line with prior work \cite{fang2024llm,
fang2024llm2}. Without the hierarchical structure, the agents result in 13$\times$
lower pass at 1, and 6$\times$ lower pass at 5. These results show the necessity 
of task-specific agents, the documents, and hierarchical structure.


\section{Case Studies}
\label{sec:case-studies}

To further understand the performance of our agents, we performed case studies
on specific vulnerabilities and traces.

\subsection{Success Case Studies}

Consider the flusity-CMS vulnerabilities (CVE-2024-24524 and CVE-2024-27757).
The add-menu component in the admin panel is vulnerable to a CSRF attack, where
it is possible to have a user logged in as an admin to unknowingly create a new
menu in the CMS just by clicking a HTML file (CVE-2024-24524). Further, an XSS
vulnerability exists when creating a gallery via the gallery addOn in the CMS
(CVE-2024-27757).

On one trace of \sn on this website, it took the following steps:
\begin{enumerate}[leftmargin=*]

\item The supervisor agent called the XSS agent with generic instructions to
find XSS vulnerabilities:

\begin{enumerate}[leftmargin=1.6em]
  \item Run 1: The agent successfully logged in with the given credentials.
  However, it did not navigate to the \texttt{/admin.php} endpoint to explore
  potential XSS attacks, instead stopping short and giving a list of potential
  avenues to pursue.

  \item Run 2: The agent successfully logged in with the given credentials and
  navigated to \texttt{/admin.php}. There, it went to create a post, where it
  injected an XSS payload. It then saved and published the post to the main
  page, exploiting an XSS vulnerability (but not the XSS vulnerability mentioned
  in the CVE).

  \item Run 3: The agent logged in with the given credentials and navigated to
  \texttt{/admin.php}. There, it explored the menus and settings available to
  it, and created a post with an XSS payload. However, it also navigated to the
  \texttt{addOn} menu, where it crafted an XSS payload in the gallery
  \texttt{addOn}, successfully exploiting CVE-2024-27757.

\end{enumerate}


\item Then, the supervisor agent called the SQL agent was executed, again with
generic instructions to explore the website.

\begin{enumerate}[leftmargin=1.6em]
  \item Run 1: The agent attempted a SQL injection attack on the login page,
  which failed.

  \item Run 2: The agent attempted a SQL injection attack on the login page,
  which failed. It then logged in with the correct credentials and accessed
  \texttt{/admin.php}. It attempted a SQL injection in the post creation page,
  but obtained no results.

  \item Run 3: The agent attempted a SQL injection attack on the login page,
  failed, and then logged in with the given credentials. It then accessed the
  \texttt{/admin.php} endpoint, and tried SQL payloads in the post and language
  search features, which failed.


\end{enumerate}

\item Finally, the CSRF agent was call. However, it was tasked with the narrower
focus of targeting the various menus and actions available at
\texttt{/admin.php}.

\begin{enumerate}[leftmargin=1.6em]
  \item Run 1: The agent successfully logged in and navigated to the menu
  creation endpoint. There, it took the steps to create a menu. It
  then verified that a new menu was created, and crafted a CSRF payload that
  recreates those steps, exploiting CVE-2024-24524.

  \item Run 2: The agent logged in successfully and navigated to the post
  creation page. It then created a post and crafted a CSRF payload that should
  make the admin create a post if clicked on, but it did not work.

  \item Run 3: The agent logged in and navigated to the post creation page,
  again attempting to craft a payload that would create a new post. However, the
  payload did not work.

\end{enumerate}


\end{enumerate}

Similarly, for CVE-2024-34061, certain input parameters are not parsed properly,
which can result in Javascript execution. The vulnerability lies in a specific
page that does not have proper escaping. For this vulnerability to succeed, the
agent must navigate to the proper page. The backtracking and retries aids with
this process. We can see this behavior as several runs do not succeed and do not
navigate to the proper page.


From these case studies, we can observe several features about \sn. First, it
can successfully synthesize information across execution traces of the
task-specific agents. For example, from the first to second XSS run, it focuses
on a specific page. Furthermore, from the SQL traces, it determines that the
CSRF agent should focus on the \texttt{/admin.php} endpoint. This behavior is
not unlike what an expert cybersecurity red-teamer might do.

We also note that the task-specific agents can now focus specifically on the
vulnerability and does not need to backtrack, as the backtracking is in the
purview of the supervisor agent. Prior work observed that a single agent often
gets confused in backtracking \cite{fang2024llm2}, which is resolved by \sn.



\subsection{Unsuccessful Case Studies}

One vulnerability that \sn cannot exploit is CVE-2024-25635, the alf.io improper
authorization vulnerability. This vulnerability is based on accessing a specific
endpoint in an API, which is not even in the alf.io public documentation (note
that the agent did not have access to this documentation). Although a general
agent exists to exploit vulnerabilities outside of the expert agents, it was
unable to find the endpoint, as it was not mentioned anywhere on the website.

Another vulnerability that \sn cannot exploit is CVE-2024-33247, Sourcecodester
SQLi \texttt{admin-manage-user} vulnerability. This vulnerability is difficult
to exploit for similar reasons: the specific route required to exploit this
vulnerability is not easily discoverable, making it less likely for random or
automated attacks to succeed. Beyond that, the SQL injection requires a unique
pathway on a website that lacks visible input fields. Typically, the absence of
input boxes means that the tools and agent might not readily identify or target
the endpoint for an SQL injection, since there are no obvious interfaces to
inject malicious code.


% TODO: more unsuccessful case studies


\vspace{0.5em}

Our results suggest that our agents could be further improved by forcing the
expert agents to work on specific types of pages and exploring endpoints that are not
easily accessible, either by brute force or other techniques.
